\uuid{yMKZ}
\titre{Quotient/Composé/DL usuels}
\auteur{hanani ; mzari}
\organisation{exo7}
\licence{CC BY-NC-SA 4.0}
\datecreate{2020-01-20}
\niveau{L1}
\module{Analyse}
\chapitre{Développement limité}
\sousChapitre{Calculs}
\difficulte{3}
\theme{Opérations sur les DL, Quotient, Composé, DL usuels}
\exo7qcmid{lille-2.developpements-limites.023}

\contenu{
\texte{
Soit $a\in \Rr$ et $f$ une fonction. On dit que $f(x)$ est un $o\left((x-a)^n\right)$ au voisinage de $a$ et on écrit $f(x)=o((x-a)^n)$ si $\displaystyle \lim _{x\to a}\frac{f(x)}{(x-a)^n}=0$. On note par $DL_n(a)f(x)$ le développement limité de $f$ à l'ordre $n$ en $a$. On note aussi $DL_n(a^+)f(x)$ (resp. $DL_n(a^-)f(x)$) le développement limité de $f$ à l'ordre $n$ à droite en $a$ (resp. à gauche en $a$).
}

\begin{question}
\qclassification{125.02}
\qtags{motcle=Quotient/Composé/DL usuels}
Parmi les assertions suivantes, cocher celles qui sont vraies :
\begin{answers}
\bad{$\displaystyle DL_2(0)\frac{\ln (1+x)}{x-x^2}=1+\frac{x}{2}+\frac{x^2}{2}+o(x^2)$.}
\bad{$\displaystyle DL_2(0)\frac{\mathrm{e}^x-1}{\ln (1+x)}=1+x+\frac{x^2}{2}+o(x^2)$.}
\good{$\displaystyle DL_4(0)\ln (1+x\sin x)=x^2-\frac{2x^4}{3}+o(x^4)$.}
\good{$\displaystyle DL_4(0)\arcsin \left(\ln (1+x^2)\right)=x^2-\frac{x^4}{2}+o(x^4)$.}
\end{answers}
\begin{explanations}
Le division suivant les puissances croissances du $\displaystyle DL_3(0)\ln (1+x)$ par $x-x^2$ donne
$$DL_2(0)\frac{\ln (1+x)}{x-x^2}=1+\frac{x}{2}+\frac{5x^2}{6}+o(x^2).$$
Le division suivant les puissances croissances du $\displaystyle DL_3(0)(\mathrm{e}^x-1)$ par le $\displaystyle DL_3(0)\ln (1+x)$ donne
$$DL_2(0)\frac{\mathrm{e}^x-1}{\ln (1+x)}=1+x+\frac{x^2}{3}+o(x^2).$$
Ensuite, avec $\displaystyle u=x\sin x =x^2-\frac{x^4}{6}+o(x^4)$ et $\ln (1+u)=u-u^2/2+o(u^2)$, on obtient :
$$DL_4(0)\ln (1+x\sin x)=x^2-\frac{2x^4}{3}+o(x^4).$$
Enfin, avec $\displaystyle u=\ln (1+x^2)=x^2-\frac{x^4}{2}+o(x^4)$ et $\arcsin u=u+\frac{u^3}{6}+o(u^4)$, on obtient :
$$DL_4(0)\arcsin \left(\ln (1+x^2)\right)=x^2-\frac{x^4}{2}+o(x^4).$$
\end{explanations}
\end{question}
}
